\subsection{An exact finite-depth dyadic obstruction from the Prym factor}

\begin{statement}{A.158}{Proved}
{Finite Selmer images, explicit valuation thresholds, and the actual quotient}
Retain \(K=\Q(\sqrt2)\), \(K_2=\Q_2(\sqrt2)\),
\(\pi=\sqrt2\), \(E_b\), \(\overline D\), and \(q_+\) from
A.150--A.154, with \(v_\pi(\pi)=1\) and \(T_0=(0,0)\).
For every \(n\geq1\), the transition induced by multiplication
by two on the coefficient modules has the exact image
\[
 \operatorname{im}\left(
 \operatorname{Sel}^{2^{n+1}}(E_b/K)
 \longrightarrow\operatorname{Sel}^{2^n}(E_b/K)\right)
       =\{\kappa_{2^n}(O),\kappa_{2^n}(T_0)\}.
\tag{A.158.1}
\]
The two classes in this set are distinct.
Its dyadic inverse image under the ordinary point Kummer map is
exactly
\[
 Z_n=[2^n]E_b(K_2)\ \cup\
                  \bigl(T_0+[2^n]E_b(K_2)\bigr).
\tag{A.158.2}
\]
For a non-torsion point \(R=(x,y)\in Z_n\), put
\(r_x=v_\pi(x)\).  Then
\[
                 r_x\leq2-4n\qquad\text{or}\qquad r_x\geq4n.
\tag{A.158.3}
\]
Thus the literal integers \(3-4n,\ldots,4n-1\) are excluded.

On the displayed harmonic curve, the actual scaled Prym coordinate is
\[
 x(q_+(T,W))
     =\frac{573M^2}{191^2}\frac{W^2}{T},
 \qquad
 r(T,W)=2v_\pi(W)-v_\pi(T),\qquad M=325596479663.
\tag{A.158.4}
\]
For every ordinary dyadic point of \(D^\circ\), this valuation
is finite.  Its image under the actual multiplication quotient
at exponent \(2^{n+1}\) cannot be the dyadic localization of
a global Selmer class whenever
\[
 3-4n\leq r(T,W)\leq4n-1.
\tag{A.158.5}
\]
In particular the explicitly chosen integer
\[
 n(T,W)=1+\left\lfloor
                    \frac{\max\{r(T,W),\,2-r(T,W)\}}4
                         \right\rfloor
\tag{A.158.6}
\]
gives such an obstruction at exponent \(2^{n(T,W)+1}\).

These are finite geometric quotient tests.
Over \(K\), compactification and \(q_+\) give the literal
surjection
\[
 \pi_1(D^\circ_K)\longrightarrow
              E_b[2^{n+1}]\rtimes G_K,
 \qquad |E_b[2^{n+1}]|=2^{2n+2}.
\tag{A.158.7}
\]
For the actual degree-eight harmonic cover \(D^\circ\to U\),
let \(\Pi=\pi_1(U_{\overline\Q})\),
\(H=\pi_1(D^\circ_{\overline\Q})\), and let \(I\) be the
closed normal subgroup of \(H\) generated by its four
boundary inertia groups.  Put
\[
 N_m=\ker\left(H\longrightarrow
              (H/I)^{\rm ab}/2^m(H/I)^{\rm ab}\right).
\tag{A.158.8}
\]
The actual arithmetic quotient of the tripod with geometric
group \(Q_m=\Pi/N_m\) has
\[
 |Q_m|=2^{6m+3},\qquad
 \operatorname{cl}(Q_m)=4m+2,\qquad \Pi''\subseteq N_m.
\tag{A.158.9}
\]
After passage to its actual \(D^\circ\) subgroup and restriction
to \(G_K\), \(Q_{n+1}\) maps to (A.158.7).
Consequently every locally point-or-cusp section of this
\(D^\circ\) subgroup has its dyadic ordinary centre in the
set (A.158.3); (A.158.5) excludes that finite local image.
The comparison uses the actual subgroup of the harmonic
cover, with its specified arithmetic twist.
In particular the test at \(m=n+1\) is a quotient of the
actual class-\((4n+6)\) geometric quotient, of order
\(2^{6n+9}\).
\end{statement}

\begin{proof}
\emph{The exact transition on Selmer groups.}
The finite multiplication sequences give a commutative diagram
whose rows are
\[
 0\longrightarrow E_b(K)/2^jE_b(K)
 \longrightarrow\operatorname{Sel}^{2^j}(E_b/K)
 \longrightarrow\Sha(E_b/K)[2^j]\longrightarrow0.
\tag{A.158.10}
\]
For \(j=n+1\to n\), the left map is the quotient map and
the right map is multiplication by two; both follow from
the actual coefficient map \(E_b[2^{n+1}]\xrightarrow{[2]}
E_b[2^n]\).
A.154 proves
\(E_b(K)=\{O,T_0\}\) and
\(\Sha(E_b/K)[2^\infty]=(\Z/2)^2\).
The right map is therefore zero.  Every image lies in the
left Kummer subgroup, and each of its two classes is the
image of the Kummer class of the same rational point at
level \(n+1\).  This proves equality (A.158.1), including
distinctness because \(T_0\) has no rational half.

For local points, the finite Kummer exact sequence has
kernel \([2^n]E_b(K_2)\).  Its inverse image of the two
classes in (A.158.1) is therefore exactly (A.158.2).
A.152 shows that every nonzero point of
\([2^n]E_b(K_2)\) has \(x\)-valuation at most \(2-4n\).
The actual group-law formula for translation by \(T_0\) is
\[
 (x,y)-T_0=(x,y)+T_0
       =\left(\frac b{x},-\frac{by}{x^2}\right),
 \qquad v_\pi(b)=2.
\tag{A.158.11}
\]
For a non-torsion \(R\), both \(R\) and \(R-T_0\) are
nonzero.  Applying A.152 to the relevant member of
(A.158.2) gives either \(r_x\leq2-4n\) or
\(2-r_x\leq2-4n\), proving (A.158.3).

\emph{The coordinate and finite depth.}
The positive-sign quotient in A.150 has unscaled coordinate
\(aW^2/T\) and integral scaling \(x\mapsto382^2x\).
Using
\[
       a=573\left(\frac{M}{2\cdot191^2}\right)^2
 \quad\text{gives}\quad
       382^2a=\frac{573M^2}{191^2}.
\tag{A.158.12}
\]
This coefficient is a dyadic unit.  Ordinary points have
\(T,W\ne0\), and A.152 identifies the complete proper
inverse image of \(\{O,T_0\}\) with the four boundary
points.  Thus (A.158.4) and (A.158.5) follow.
The strict inequalities
\(4n>r\) and \(4n>2-r\) are precisely the negation of
(A.158.3).  Their least positive integer solution is
(A.158.6).

\emph{The multiplication quotient and the harmonic comparison.}
The covers \([2^m]:E_b\to E_b\), pointed by the origin,
give the arithmetic quotient \(E_b[2^m]\rtimes G_K\).
The open immersion \(D^\circ_K\to\overline D_K\) kills
the four boundary inertia groups.  The proper morphism
\(q_+\) gives the next map, and it sends the rational
boundary point \(P_0\) to \(O\).
Its geometric fundamental-group map is onto: its image
is closed, and every proper closed subgroup of a profinite
group lies in a proper open subgroup, as is seen by
separating an element outside it in a finite quotient.
Thus a proper image, including one of infinite index,
would give a nontrivial connected finite
etale intermediate cover of \(E_b\) through which the
degree-two map \(q_+\) factors.  Its degree would be two,
so \(q_+\) itself would be etale and the source would have
genus one, contradicting its actual genus three.
Equivalently, this is the connected-cover correspondence
applied to the image of the proper fundamental group.
The surjection remains after the open immersion and gives
(A.158.7).  Rational tangential specialization at \(P_0\)
maps to the origin section, so the displayed semidirect
action is the actual Galois action on elliptic torsion.

The degree-eight cover has normal geometric subgroup \(H\)
and quotient \(C_8\), as proved by its literal equation
in A.139.  Arithmetic conjugation permutes the four
boundary inertia groups, hence preserves \(I\) and \(N_m\).
The proper genus-three geometric group has abelianization
\(\widehat\Z^{\,6}\).  Thus
\(H/N_m=(\Z/2^m)^6\), proving the order in (A.158.9).
Since \(\Pi/H\) is abelian, \(\Pi'\subset H\); since
\(H/N_m\) is abelian, \(\Pi''\subset[H,H]\subset N_m\).
These are the same subgroup and boundary maps as A.157.
For clarity, its integral presentation gives the actual
semidirect product
\[
 Q_m=R_m\rtimes C_8,\qquad
 R_m=(\Z/2^m)[V]/(V^6+V^4+V^2+1),
 \qquad vuv^{-1}=Vu.
\tag{A.158.13}
\]
Here \(v\) generates \(C_8\), \(u\) is the constant-one
module element, and the original tripod generators are
\(x=uv^2,\ y=v\).  This follows from the Schreier
generators \(u_j=v^juv^{-j}\) and \(z=v^8\): proper
boundary specialization kills \(z\), the even sum of
the \(u_j\), and the odd sum.  The two sums are the
polynomial \(1+V^2+V^4+V^6\) and its \(V\)-multiple.

We prove the exact class, retaining the integral module.
Put \(\delta=V-1\), \(f=V^2+1\), \(g=V^4+1\), and
\(R=\Z_2[V]/(fg)\).  The literal identity
\(g=(V^2-1)f+2\), together with coprimality of \(f,g\)
over \(\Q_2[V]\), gives
\[
 R=
 \Z_2[i]\mathbin{\times}_{\,\F_2[\delta]/(\delta^2)}
                   \Z_2[\zeta_8].
\tag{A.158.14}
\]
Indeed \((f)+(g)=(f,2)\), and \((f)\cap(g)=(fg)\)
in the unique factorization domain \(\Z_2[V]\).
The two translated polynomials \(f(1+\delta)\) and
\(g(1+\delta)\) are Eisenstein at two.  Thus \(\delta\)
is a uniformizer in both factors, with respective
valuations of two equal to two and four.
The pair \(\delta^{4m+2}/2^m\) has respective
uniformizer valuations \(2m+2\) and two.  Both
reductions in the common quotient of (A.158.14) are
zero, so this pair belongs to \(R\).  In contrast,
\(\delta^{4m+1}/2^m\) has valuations \(2m+1\) and one:
its first reduction is zero and its second is nonzero.
It does not belong to \(R\).  Consequently
\[
 \delta^{4m+2}R\subseteq2^mR,\qquad
                      \delta^{4m+1}\notin2^mR.
\tag{A.158.15}
\]
The semidirect product (A.158.13) has abelian kernel
and cyclic quotient, with
\(\Gamma_jQ_m=\delta^{j-1}R_m\) for every \(j\geq2\).
These are actual commutators with \(v\), and the
orbit of \(u\) generates the whole module.  Equation
(A.158.15) proves exact class \(4m+2\).
The map to \(E_b[2^m]\) factors through proper abelian
homology modulo \(2^m\), so it kills \(N_m\).
This proves the stated map from the actual \(D^\circ\)
subgroup of \(Q_m\), after restriction to \(G_K\).

A section locally represented by an ordinary \(D^\circ\)
point maps to the Kummer class of its image under \(q_+\).
A local cusp maps, after compactification, to the proper
boundary point and then to its elliptic Kummer class.
Thus these sections produce elements of
\(\operatorname{Sel}^{2^m}(E_b/K)\) through the actual
maps, with their local images preserved.  Applying
(A.158.1) for \(m=n+1\) proves the finite quotient assertion.

\emph{The behavior near the boundary.}
For each fixed \(n\), the map \([2^{n+1}]\) on \(E_b(K_2)\)
has nonzero derivative at the origin.  The inverse
function form of Hensel's lemma therefore makes its
image contain a neighborhood of the origin.
The map \(q_+\) is etale at each of the four boundary
points by its exact fibres in A.152.
Its local inverse on a sufficiently small neighborhood
therefore provides ordinary points arbitrarily close
to each boundary point whose finite Kummer class at
exponent \(2^{n+1}\) is exactly the global class of
the corresponding torsion point.  In particular their
images belong to the corresponding coset in (A.158.2).
This proves directly that the depth in (A.158.6) grows
near the boundary.  For the displayed elliptic finite
Kummer test, no single finite level excludes all ordinary
dyadic points.
\end{proof}

\paragraph{One fixed dyadic centre for the entire family of A.145.}
For the frozen dyadic point used in A.134 and A.145,
\[
 T_2=\frac{31}{2},\qquad
 X_0\equiv5,\quad X_1\equiv3,\quad X_h\equiv5\pmod8,
 \qquad W=X_0X_1X_h.
\tag{A.158.16}
\]
Thus \(W\) is a unit, \(v_\pi(T_2)=-2\), and
\[
                 r(T_2,W)=2,\qquad
                       t_2=4T_2(1-T_2)=-899.
\tag{A.158.17}
\]
Every \(\Q_2\)-point of the actual degree-eight cover
above \(t=-899\) has this same \(T\)-coordinate.
The other possible coordinate is \(1-T_2\).
The polynomial \(F(T)=T(T-1)(T-\tfrac12)\) satisfies
\(F(1-T)=-F(T)\), so that alternative would require
\((W'/W)^4=-1\) in \(\Q_2\).
Such an element would be a unit, whereas every odd
fourth power is one modulo sixteen and \(-1\) is
fifteen modulo sixteen.  This proves the exclusion
of the alternative sheet without a choice of lift.

Consequently there is no section of the arithmetic
quotient \(E_U/N_2\) with combined mod-eight class
\([64a^2]\), with dyadic image conjugate to the ordinary
point image at \(t=-899\), and with images of actual
ordinary or cuspidal local sections at every other
place, including the real place.  Indeed the exact
fixed-sheet comparison of A.157 lifts such a section
to the \(D^\circ\) subgroup; its dyadic lift has the
valuation two in (A.158.17).  Taking \(n=1\) in
(A.158.5) contradicts the finite Selmer image.
Here \(Q_{m=2}\) denotes the finite geometric group,
of order \(2^{15}=32768\) and exact class ten;
\(\Q_2\) throughout denotes the local field.
The other local centres are unrestricted.  This one
finite quotient therefore applies to every member
of the family of A.145 with its frozen dyadic image.

The certificate
\path{code/harmonic_prym_finite_depth_certificate.py}
checks the exact scaling and torsion translation, all threshold
and iteration identities on the displayed finite ranges, and
literal Hensel points approaching each of the four boundary
points.  These finite checks do not replace the all-depth proof.
